% GENERATED FILE. Edit canonical lessons, then run npm run generate:books.
% canonical-source-hash: fnv1a64:4ad836ecf8536459
% canonical-lessons: SA-C518-tulika, SA-C518-rangah, SA-C518-rekha, SA-C518-alikhati, SA-C518-pratikrtih

\chapter{Colours on the Wall}
\label{ch:sa-colours-on-the-wall}
\hlchaptermodality{518}

\begin{chapteropening}
\textbf{By the end of this chapter:} \emph{I can talk about brushes, lines and paint.}

Complete the last lesson of chapter 518: I can talk about brushes, lines and paint.
\end{chapteropening}

\section[\texorpdfstring{\sk{तूलिका} (tūlikā)}{तूलिका (tūlikā)}]{\sk{तूलिका} (tūlikā) — a paintbrush}

\label{lesson:SA-C518-tulika}



\begin{quote}
Before the new one: say the Sanskrit for a marble, then the Sanskrit for a doll.
\end{quote}

\subsection*{You'll want to know: \sk{तूलिका}}
\textbf{\sk{तूलिका}} — \emph{tūlikā} — \textquotedblleft{}a paintbrush\textquotedblright{}.

\begin{quote}
\sk{चित्रकारः} \sk{नूतनां} \sk{तूलिकां} \sk{क्रीणाति।}
\end{quote}

\begin{quote}
\emph{citrakāraḥ nūtanāṃ tūlikāṃ krīṇāti.}
\end{quote}

\begin{quote}
The painter buys a new brush.
\end{quote}

A feminine noun.

\begin{cousinweb}
\textbf{\sk{तूलिका}} (\emph{tūlikā}) is the word MW gives as the same as \emph{tūli}, a painter's brush; MW gives no further derivation, so its origin is uncertain (MW).
\end{cousinweb}

\begin{grammarlens}[title={What you've built}]
One more word for describing what you lived and felt.
\end{grammarlens}

\subsection*{Your turn}
\begin{itemize}
\raggedright
  \item \emph{Say it:} \emph{tūlikā}
  \item \emph{Say it:} \emph{tūlikā}, once more
  \item \emph{Say it:} \emph{citrakāraḥ nūtanāṃ tūlikāṃ krīṇāti.} — The painter buys a new brush.
  \item \emph{Recall:} say the Sanskrit for a marble, then the Sanskrit for a doll, then say \emph{tūlikā} again
\end{itemize}

\begin{tcolorbox}[breakable,colback=teal!4,colframe=teal!35!black,title={Before you move on}]
What does \sk{तूलिका} mean? (\textquotedblleft{}A paintbrush\textquotedblright{}.) Say it once more.
\end{tcolorbox}
\section[\texorpdfstring{\sk{रंगः} (raṅgaḥ)}{रंगः (raṅgaḥ)}]{\sk{रंगः} (raṅgaḥ) — paint; a dye}

\label{lesson:SA-C518-rangah}



\begin{quote}
Before the new one: say the Sanskrit for a doll, then the Sanskrit for a paintbrush.
\end{quote}

\subsection*{You'll want to know: \sk{रंगः}}
\textbf{\sk{रंगः}} — \emph{raṅgaḥ} — \textquotedblleft{}paint; a dye\textquotedblright{}.

\begin{quote}
\sk{चित्रकारः} \sk{रंगं} \sk{जलेन} \sk{संयोजयति।}
\end{quote}

\begin{quote}
\emph{citrakāraḥ raṅgaṃ jalena saṃyojayati.}
\end{quote}

\begin{quote}
The painter mixes the paint with water.
\end{quote}

A masculine noun.

\begin{cousinweb}
\textbf{\sk{रंगः}} (\emph{raṅgaḥ}) comes from the root √rañj, \textquotedblleft{}to be coloured\textquotedblright{}; MW gives \emph{raṅga} as colour, paint, dye, and also as a theatre, a stage (MW).
\end{cousinweb}

\begin{grammarlens}[title={What you've built}]
One more word for describing what you lived and felt.
\end{grammarlens}

\subsection*{Your turn}
\begin{itemize}
\raggedright
  \item \emph{Say it:} \emph{raṅgaḥ}
  \item \emph{Say it:} \emph{raṅgaḥ}, once more
  \item \emph{Say it:} \emph{citrakāraḥ raṅgaṃ jalena saṃyojayati.} — The painter mixes the paint with water.
  \item \emph{Recall:} say the Sanskrit for a doll, then the Sanskrit for a paintbrush, then say \emph{raṅgaḥ} again
\end{itemize}

\begin{tcolorbox}[breakable,colback=teal!4,colframe=teal!35!black,title={Before you move on}]
What does \sk{रंगः} mean? (\textquotedblleft{}Paint; a dye\textquotedblright{}.) Say it once more.
\end{tcolorbox}
\section[\texorpdfstring{\sk{रेखा} (rekhā)}{रेखा (rekhā)}]{\sk{रेखा} (rekhā) — a line}

\label{lesson:SA-C518-rekha}



\begin{quote}
Before the new one: say the Sanskrit for a paintbrush, then the Sanskrit for paint.
\end{quote}

\subsection*{You'll want to know: \sk{रेखा}}
\textbf{\sk{रेखा}} — \emph{rekhā} — \textquotedblleft{}a line\textquotedblright{}.

\begin{quote}
\sk{बालः} \sk{ऋजुं} \sk{रेखां} \sk{लिखति।}
\end{quote}

\begin{quote}
\emph{bālaḥ ṛjuṃ rekhāṃ likhati.}
\end{quote}

\begin{quote}
The child draws a straight line.
\end{quote}

A feminine noun.

\begin{cousinweb}
\textbf{\sk{रेखा}} (\emph{rekhā}) probably comes from the root √likh (older \emph{rikh}), \textquotedblleft{}to scratch\textquotedblright{}, beside which MW sets it; MW gives \emph{rekhā} as a scratch, a streak, a line, a row (MW).
\end{cousinweb}

\begin{grammarlens}[title={What you've built}]
One more word for describing what you lived and felt.
\end{grammarlens}

\subsection*{Your turn}
\begin{itemize}
\raggedright
  \item \emph{Say it:} \emph{rekhā}
  \item \emph{Say it:} \emph{rekhā}, once more
  \item \emph{Say it:} \emph{bālaḥ ṛjuṃ rekhāṃ likhati.} — The child draws a straight line.
  \item \emph{Recall:} say the Sanskrit for a paintbrush, then the Sanskrit for paint, then say \emph{rekhā} again
\end{itemize}

\begin{tcolorbox}[breakable,colback=teal!4,colframe=teal!35!black,title={Before you move on}]
What does \sk{रेखा} mean? (\textquotedblleft{}A line\textquotedblright{}.) Say it once more.
\end{tcolorbox}
\section[\texorpdfstring{\sk{आलिखति} (ālikhati)}{आलिखति (ālikhati)}]{\sk{आलिखति} (ālikhati) — paints}

\label{lesson:SA-C518-alikhati}



\begin{quote}
Before the new one: say the Sanskrit for paint, then the Sanskrit for a line.
\end{quote}

\subsection*{You'll want to know: \sk{आलिखति}}
\textbf{\sk{आलिखति}} — \emph{ālikhati} — \textquotedblleft{}paints\textquotedblright{}.

\begin{quote}
\sk{मीरा} \sk{पुष्पाणि} \sk{आलिखति।}
\end{quote}

\begin{quote}
\emph{mīrā puṣpāṇi ālikhati.}
\end{quote}

\begin{quote}
Mira paints flowers.
\end{quote}

Past (she …): \sk{सा} \sk{मयूरम्} \sk{आलिखितवती}, \emph{sā mayūram ālikhitavatī}.

\begin{cousinweb}
\textbf{\sk{आलिखति}} (\emph{ālikhati}) is built on \emph{ā} and the root √likh, \textquotedblleft{}to scratch, to write\textquotedblright{}, which also gives \emph{likhati}; MW gives \emph{ālikh} as to mark, to draw, to delineate, to paint (MW).
\end{cousinweb}

\begin{grammarlens}[title={What you've built}]
One more word for describing what you lived and felt.
\end{grammarlens}

\subsection*{Your turn}
\begin{itemize}
\raggedright
  \item \emph{Say it:} \emph{ālikhati}
  \item \emph{Say it:} \emph{ālikhati}, once more
  \item \emph{Say it:} \emph{mīrā puṣpāṇi ālikhati.} — Mira paints flowers.
  \item \emph{Recall:} say the Sanskrit for paint, then the Sanskrit for a line, then say \emph{ālikhati} again
\end{itemize}

\begin{tcolorbox}[breakable,colback=teal!4,colframe=teal!35!black,title={Before you move on}]
What does \sk{आलिखति} mean? (\textquotedblleft{}Paints\textquotedblright{}.) Say it once more.
\end{tcolorbox}
\section[\texorpdfstring{\sk{प्रतिकृतिः} (pratikṛtiḥ)}{प्रतिकृतिः (pratikṛtiḥ)}]{\sk{प्रतिकृतिः} (pratikṛtiḥ) — a portrait; a likeness}

\label{lesson:SA-C518-pratikrtih}



\begin{quote}
Before the new one: say the Sanskrit for a line, then the Sanskrit for paints.
\end{quote}

\subsection*{You'll want to know: \sk{प्रतिकृतिः}}
\textbf{\sk{प्रतिकृतिः}} — \emph{pratikṛtiḥ} — \textquotedblleft{}a portrait; a likeness\textquotedblright{}.

\begin{quote}
\sk{चित्रकारः} \sk{मम} \sk{मातुः} \sk{प्रतिकृतिम्} \sk{आलिखति।}
\end{quote}

\begin{quote}
\emph{citrakāraḥ mama mātuḥ pratikṛtim ālikhati.}
\end{quote}

\begin{quote}
The painter paints my mother's portrait.
\end{quote}

A feminine noun.

\begin{cousinweb}
\textbf{\sk{प्रतिकृतिः}} (\emph{pratikṛtiḥ}) is built on \emph{prati}, \textquotedblleft{}in return, against\textquotedblright{}, and the root √kṛ, \textquotedblleft{}to make\textquotedblright{}: a thing made to answer another; MW gives \emph{pratikṛti} as an image, a likeness, a model (MW).
\end{cousinweb}

\begin{grammarlens}[title={What you've built}]
One more word for describing what you lived and felt.
\end{grammarlens}

\subsection*{Your turn}
\begin{itemize}
\raggedright
  \item \emph{Say it:} \emph{pratikṛtiḥ}
  \item \emph{Say it:} \emph{pratikṛtiḥ}, once more
  \item \emph{Say it:} \emph{citrakāraḥ mama mātuḥ pratikṛtim ālikhati.} — The painter paints my mother's portrait.
  \item \emph{Recall:} say the Sanskrit for a line, then the Sanskrit for paints, then say \emph{pratikṛtiḥ} again
\end{itemize}

\begin{tcolorbox}[breakable,colback=teal!4,colframe=teal!35!black,title={Before you move on}]
What does \sk{प्रतिकृतिः} mean? (\textquotedblleft{}A portrait; a likeness\textquotedblright{}.) Say it once more.
\end{tcolorbox}
